\begin{abstract}
Forecasts of AI-linked outcomes commonly extrapolate benchmark performance, compute, price, adoption, or economic output as though they were interchangeable. That shortcut is structurally fragile: capability is not reliable execution; reliable execution is not verified completion; verified completion is not authorized deployment; deployment is not adoption; and adoption is not realized economic or mission output. The problem becomes more severe when capability, cost, task duration, automation, parallelism, reliability, demand, and institutional constraints evolve at different rates, and when the law governing those rates may itself change during the forecast horizon. We introduce an acceleration-aware, constraint-limited, regime-switching framework for forecasting realized outcomes. Its compact canonical model defines realized output as the residual-bottleneck-adjusted minimum of technical capacity and absorptive demand. A generalized realization family then separates technical, demand, and transfer bottlenecks and permits partial substitution while retaining the hard-min model as an interpretable limiting case. Technical-output growth is decomposed into capability, cost decline, task-time compression, automation, effective parallelism, and verified reliability. A task-graph layer maps model-level progress into retries, verification, human fallback, authority latency, critical paths, correlated failure, and cost per verified outcome. The reference engine fits competing regimes, generates residual-bootstrap model-averaged predictive distributions, applies downside, accelerated, and hazard-driven discontinuity transformations, calculates milestone probabilities, renders a self-contained report, and can append the frozen forecast to a separate hash-chained registry after an explicit registration step. The release supplies a complete natural-language protocol, strict machine-readable contracts, deterministic controlled experiments, source-linked METR and SWE-bench demonstrations, an operational-evidence panel, a prospective validation design, and reproducible release infrastructure. The paper does not claim that each constituent equation is novel or that the demonstrations prospectively validate future AI trajectories. Its contribution is an integrated, auditable, falsifiable bridge from fresh evidence to verified, authorized, and realized outcomes in a world whose rate of change may itself change.
\end{abstract}

\begin{keybox}
\textbf{Research contribution.} The framework makes the following chain explicit and testable:
\[
\begin{gathered}
\text{capability}\neq\text{reliable execution}\neq\text{verified completion},\\[0.15em]
\text{authorized deployment}\neq\text{adoption}\neq\text{realized value},\\[0.2em]
\text{forecast}\rightarrow\text{dated evidence}\rightarrow\text{probabilistic scenarios}\rightarrow\text{action}\rightarrow\text{prospective score}.
\end{gathered}
\]
It is designed to prevent two symmetric failures: converting technical gains directly into deployment or revenue, and preserving historical execution timelines after automation, parallelism, verification cost, or authority latency have materially changed.
\end{keybox}

\textbf{Keywords:} technological forecasting; artificial intelligence; endogenous acceleration; regime switching; task graphs; reliability; adoption; bottlenecks; probabilistic forecasting; hindcasting; prospective validation; decision-making under deep uncertainty.

\textbf{Status and disclosure.} Version 2.0.0 is a versioned research preprint, executable protocol, and reference forecasting engine. It is not an audit, certification, assurance engagement, investment recommendation, legal opinion, or promise of forecast accuracy. The principal author directed the research design, assumptions, editorial decisions, and release. Generative AI systems assisted with research, drafting, coding, testing, and document production. The author remains responsible for the published content. No DOI, public repository state, or external preprint identifier is claimed until one is actually issued.

\textbf{License.} Software: AGPL-3.0-or-later. Paper, prompt, schemas, documentation, and original research artifacts: CC BY-NC-SA 4.0, subject to the repository's license map. Trademarks reserved. See \texttt{LICENSE.md}.
